\subsection{Subrepresentations of the regular representation of a commutative group}

In this section, G is a locally compact commutative group and \(\mu\) a Haar measure on G. We denote by \(\hat{G}\) the dual group of G (def. 2 of II, p. 201) and by \(\hat{\mu}\) the dual Haar measure of \(\mu\) onto \(\hat{G}\) (def. 4 of II, p. 214). The notions of measurability will always be relative to \(\mu\) and \(\hat{\mu}\) .

We propose to determine all subrepresentations of the left regular representation \(\gamma_{G}\) of G in \(L^{2}(G,\mu)\) . Since G is commutative, we also have \(\delta_{\mathrm{G}}(g)=\gamma_{\mathrm{G}}(g^{-1})\) , so these subrepresentations are also the subrepresentations of the right regular representation.

For every measurable subset M of \(\widehat{\mathbf{G}}\) , we denote \(\mathrm{E_M}\) the set of \(f \in \mathrm{L}^2 (\mathrm{G},\mu)\) such that the Fourier transform \(\mathcal{F}_{\mathrm{G}}(f)\) is zero almost everywhere on M (cf. no. 3 of II, p. 210). This is the kernel of the continuous linear map \(f \mapsto \varphi_{\mathrm{M}}\mathcal{F}_{\mathrm{G}}(f)\) of \(\mathrm{L}^2 (\mathrm{G},\mu)\) in \(\mathrm{L}^2 (\widehat{\mathrm{G}},\widehat{\mu})\) , where \(\varphi_{\mathrm{M}}\) denotes the characteristic function of M (cf. th. 1 of II, p. 215), and it is therefore a closed subspace of \(\mathrm{L}^2 (\mathrm{G},\mu)\) .

We say that measurable subsets M and N of \(\widehat{G}\) are equal up to a locally negligible set if \((\mathrm{M} \cup \mathrm{N}) - (\mathrm{M} \cap \mathrm{N})\) is locally negligible. Equivalently, measurable subsets M and N are equal up to a locally negligible set if and only if the characteristic functions of M and N are equal in \(\mathrm{L}^{\infty}(\widehat{\mathrm{G}}, \widehat{\mu})\) .

PROPOSITION 10. --- a) Let M be a measurable subset of \(\widehat{G}\) . The space \(\mathrm{E}_{\mathrm{M}}\) is a subrepresentation of the representation \(\gamma_{\mathrm{G}}\) ;

\begin{enumerate}
\def\labelenumi{\alph{enumi})}
\setcounter{enumi}{1}
\item
  Let M and N be measurable subsets of \(\widehat{\mathbf{G}}\) . We have \(\mathrm{E_M} = \mathrm{E_N}\) if and only if M and N are equal up to a locally negligible set;
\item
  Any subrepresentation of \(\gamma_{\mathrm{G}}\) is of the form \(\mathrm{E}_{\mathrm{M}}\) for a measurable subset M of \(\widehat{\mathrm{G}}\) .
\end{enumerate}

Let \(\eta\) the canonical map from G to \(\widehat{\mathbf{G}}\) (cf. II, p. 216, remark 1). Since \(\mathrm{E_M}\) is closed, the assertion \(a\) ) follows from the formulas

\[
\boldsymbol {\gamma} _ {\mathrm{G}} (x) (f) = \varepsilon_ {x} * f, \quad \mathcal {F} _ {\mathrm{G}} (\varepsilon_ {x} * f) = \eta (x) \mathcal {F} _ {\mathrm{G}} (f)
\]

for \(x\in \mathrm{G}\) and \(f\in \mathrm{L}^2 (\mathrm{G},\mu)\) .

Let M and N be measurable subsets equal up to a locally negligible set. Since \(\varphi_{M}\) is then equal to \(\varphi_{N}\) in \(L^{\infty}(\widehat{G},\widehat{\mu})\) , this condition implies that \(E_{M}=E_{N}\) .

Conversely, suppose that M and N are not equal up to a locally negligible set. Possibly after interchanging the roles of M and N, there then exists a compact subset K such that the set \(L = K \cap (M - (M \cap N))\) is not negligible. Let \(\varphi_{\mathrm{L}} \in \mathrm{L}^{2}(\widehat{\mathrm{G}}, \widehat{\mu})\) the class of the characteristic function of L, and set \(f = \overline{\mathcal{F}}_{\widehat{\mathrm{G}}}(\varphi_{\mathrm{L}}) \in \mathrm{L}^{2}(\mathrm{G}, \mu)\) ; we have \(\mathcal{F}_{\mathrm{G}}(f) = \varphi_{\mathrm{L}}\) (corollary of Theorem 2 in II, p. 220). Then f belongs to \(E_{N}\) since \(L \cap N\) is empty, but not in \(E_{M}\), since \(\varphi_{\mathrm{M}}\mathcal{F}(f) = \varphi_{\mathrm{M}}\varphi_{\mathrm{L}} = \varphi_{\mathrm{L}}\) . We therefore have \(E_{M} \neq E_{N}\) , which proves b).

Let now E be a subrepresentation of \(\gamma_{\mathrm{G}}\) . Let \(p_{\mathrm{E}}\) the orthogonal projector from \(\mathrm{L}^2 (\mathrm{G},\mu)\) with image E, and let \(q_{\mathrm{E}} = \mathcal{F}_{\mathrm{G}}\circ p_{\mathrm{E}}\circ \mathcal{F}_{\mathrm{G}}^{-1}\) . The projector \(p_{\mathrm{E}}\) belongs to \(\mathrm{Hom}_{\mathrm{G}}(\gamma_{\mathrm{G}},\gamma_{\mathrm{G}})\) (prop. 4 of V, p. 383), hence it commutes with \(\gamma_{\mathrm{G}}(f)\) for all \(f\in \mathrm{L}^{1}(\mathrm{G},\mu)\) (cf. V, p. 401). This means that it commutes with the endomorphisms \(\varphi \mapsto f*\varphi\) for \(f\in \mathrm{L}^{1}(\mathrm{G},\mu)\) . By lemma 4 of V, p. 407, consequently, the endomorphism \(q_{\mathrm{E}}\) of \(\mathrm{L}^2 (\widehat{\mathrm{G}},\widehat{\mu})\) commutes with the endomorphism of multiplication by \(g\) for every function \(g\in \mathcal{C}_0(\widehat{\mathrm{G}})\) belonging to the image of the Fourier transform of \(\mathrm{L}^1 (\mathrm{G},\mu)\) in \(\mathcal{C}_0(\widehat{\mathrm{G}})\) (prop. 14 of II, p. 223). Since the image of the Fourier transform is dense in \(\mathcal{C}_0(\widehat{\mathrm{G}})\) (corollary of Prop. 5 of II, p. 209), the continuity of the morphism \(g\mapsto m_g\) implies that \(q_{\mathrm{E}}\) commutes with \(m_g\) for every function \(g\in \mathcal{C}_0(\widehat{\mathrm{G}})\) .

By the corollary of Proposition 7 of IV, p. 188, there therefore exists a function \(\varphi \in \mathcal{L}^{\infty}(\widehat{\mathrm{G}},\widehat{\mu})\) such that \(q_{\mathrm{E}} = m_{\varphi}\) . Since \(q_{\mathrm{E}} = q_{\mathrm{E}}^{2}\) , we have \(m_{\varphi} = m_{\varphi}^{2} = m_{\varphi^{2}}\) , whence \(\varphi = \varphi^{2}\) in \(\mathrm{L}^{\infty}(\widehat{\mathrm{G}},\widehat{\mu})\) (prop. 5 of IV, p. 186); this means that \(\varphi\) is equal in \(\mathrm{L}^{\infty}(\widehat{\mathrm{G}},\widehat{\mu})\) to the class of the characteristic function of a measurable subset N of \(\widehat{\mathrm{G}}\) . Set \(\mathrm{M} = \widehat{\mathrm{G}} - \mathrm{N}\) . Let \(f \in \mathrm{L}^{2}(\mathrm{G},\mu)\) . We have \(f \in \mathrm{E}\) if and only if \(p_{\mathrm{E}}(f) = f\) if and only if \(\varphi \mathcal{F}_{\mathrm{G}}(f) = \mathcal{F}_{\mathrm{G}}(f)\) , which is equivalent to \(f \in \mathrm{E}_{\mathrm{M}}\) .

Note. --- Let \(\chi\) a character of G. If G is not compact, the \(\chi\) -isotypic component of the regular representation of G in L \(^{2}\) (G)

is trivial. If G is compact, the \(\chi\) -isotypic component of the regular representation of G has dimension 1, and the function \(\chi \in \mathrm{L}^{2}(\mathrm{G})\) is a basis for it.

